\chapter{Compact integer encodings}
\MARSChapterHeading{Compact integer encodings}
\label{revision:encodings}

A sequence of readings may change much less than its absolute magnitude suggests.
The values 1000, 1002, 1001, 1001 and 20000 need twenty bytes as five words, but
most of their successive differences are small. We can store those differences
in a representation that gives small values fewer bytes. Decoding must recover
the exact readings, including a sudden change, and reject an incomplete or
invalid spelling without corrupting earlier results.

The complete \texttt{delta-codec.asm} in \texttt{revision/labs/files} contains
an encoder and decoder. The example connects signed arithmetic, logical shifts,
byte addressing and the bounded interfaces from earlier chapters. Its central
invariant separates a candidate being examined from the prefix already committed.
The independent exercise validates a stream and reports its last valid reading
without creating a decoded array.

\section{Subtract in the width the difference requires}

Start each call with predecessor zero. For reading $x_i$, form
$d_i=x_i-p$, then let $p=x_i$ after accepting that reading. For the five readings
above, the differences are $1000,2,-1,0,18999$. The repeated 1001 needs a zero
difference. Adding each difference to its predecessor reverses the operation;
no approximation or quantization is involved.

Both readings fit in signed words, but their difference may not. Moving from
$-2147483648$ to 2147483647 gives 4294967295, while the reverse gives
$-4294967295$. A word subtraction would retain a wrapped pattern instead of
that exact value. The worked encoder loads with \texttt{lw}, obtaining a
sign-extended reading in a 64-bit register, then uses \texttt{dsubu} against
the full predecessor. Every possible difference fits in a signed doubleword.

A complete encoded block needs an agreed starting predecessor. This format uses
zero for every call. Splitting an array into two independently encoded blocks
therefore restarts the second block at zero. Concatenating their bytes and
decoding as one block would wrongly apply the first block's final predecessor
to the second block's first difference. Either keep block boundaries or design
an interface that explicitly carries predecessor state. Reusing a register from
a previous call is not a substitute for such an interface.

\section{Put small negative and positive values near zero}

An unsigned representation naturally gives small nonnegative values short
spellings. Map a signed difference to an unsigned integer $u$ by
\[
 u=\begin{cases}2d,&d\geq0,\\-2d-1,&d<0.\end{cases}
\]
This interleaves the two sides of zero: differences $0,-1,1,-2,2$ become
$0,1,2,3,4$. This is commonly called ZigZag mapping. Positive differences become
even integers, negative differences odd integers. It is different from copying
a negative two's-complement pattern into an unsigned interpretation, which would
place even minus one near the top of the unsigned range.

The worked implementation uses the doubleword difference in \texttt{\$t1}:
\Needspace{9\baselineskip}
\MARSCodeLabel{Map an exact signed difference}
\begin{minted}{mars-idea}
dsll $t2, $t1, 1
dsra $t3, $t1, 63
xor $t2, $t2, $t3
\end{minted}

For nonnegative $d$, the arithmetic sign shift produces zero, leaving $2d$.
For negative $d$, it produces all ones; XOR complements the doubled pattern,
which represents $-2d-1$ under these bounds. Difference $-1$ becomes one.
A logical sign shift would produce only a one bit and would not perform that
complement. Word shifts would also lose information for large differences.
The mapped value needs as many as 33 bits; a 32-bit unsigned temporary is too small.

To reverse the mapping, let $q=\lfloor u/2\rfloor$. An even $u$ gives $d=q$;
an odd $u$ gives $d=-q-1$. The decoder uses the low bit to form either a zero
mask or an all-one mask, then XORs that mask with the logical half. The sign of
the original reading is not encoded separately: it emerges only after adding
the recovered difference to the predecessor.

\section{Store seven payload bits in each byte}

Write $u$ in base 128, least significant digit first. A byte's low seven bits
hold one digit; its high bit is one when another digit follows and zero in the
last byte. The number of bytes is variable, so this is a \emph{varint}. The
zero value still needs one byte, hexadecimal \texttt{00}.

For the first difference 1000, mapped value 2000 has digits 80 and 15 because
$2000=80+128\cdot15$. The first digit becomes hexadecimal \texttt{d0}, with
its continuation bit set; the terminal digit is \texttt{0f}. This is an
ordering of seven-bit digits, not a two-byte little-endian word store.
The next differences $2,-1,0$ become bytes \texttt{04}, \texttt{01},
\texttt{00}. The last maps to 37998, giving \texttt{ee a8 02}.

\begin{figure}[htbp]
  \centering
  \begin{tabular}{r r r l}
    \toprule
    Reading & Exact difference & Unsigned mapping & Wire bytes (hexadecimal) \\
    \midrule
    1000    & 1000             & 2000             & \texttt{d0 0f}           \\
    1002    & 2                & 4                & \texttt{04}              \\
    1001    & $-1$             & 1                & \texttt{01}              \\
    1001    & 0                & 0                & \texttt{00}              \\
    20000   & 18999            & 37998            & \texttt{ee a8 02}        \\
    \bottomrule
  \end{tabular}
  \caption{The worked stream begins with predecessor zero. Each row becomes one complete varint; concatenating the rows produces eight bytes. With byte capacity seven, the two slots remaining after the five-byte committed prefix cannot hold the final three-byte candidate.}
  \label{fig:revision-delta-wire}
\end{figure}


The byte loop keeps the remaining unsigned value in \texttt{\$t2} and a full
output cursor in \texttt{\$t5}. It masks the low seven bits, logically shifts
the remaining value right by seven, and tests whether anything remains.
For a nonterminal digit it adds the high bit before storing. The final digit
is stored without that bit. The complete caller's five readings therefore
produce eight bytes, which the decoder must reconstruct exactly.

The format requires the shortest spelling. A multibyte value whose terminal
payload is zero has an unnecessary final group: \texttt{80 00} is a redundant
spelling of zero, and \texttt{81 00} a redundant spelling of one. Both are
malformed here. An internal zero digit is permitted, as in \texttt{80 01}
for unsigned 128; only the unnecessary highest group violates minimality.

At most five bytes are permitted. Four groups provide 28 payload bits, leaving
five bits for the fifth group. That byte must terminate and have payload less
than 32. A fifth byte with its high bit set is already malformed; the decoder
must not fetch a sixth byte to find out what happens. These syntax rules define
a local format, not every format that uses a base-128 integer.

\section{Reserve a complete encoding before storing it}

The encoder receives a full word pointer and reading count in
\texttt{\$a0/\$a1}, then a full byte destination and byte capacity in
\texttt{\$a2/\$a3}. Counts range from zero through 32, capacities from zero
through 160. Its input and output slices are disjoint. The function returns
committed readings in \texttt{\$v0}, committed bytes in \texttt{\$v1}, and
status in \texttt{\$a3}: zero complete or one capacity exhausted.

After computing a candidate's mapped value, the worked encoder counts how many
base-128 bytes it needs using a copy of that value. Length starts at one, even
for zero. Let $C$ be total byte capacity, $B$ committed bytes, and $L$ candidate
length. Test $L\leq C-B$ before any candidate store. The invariant $B\leq C$
makes the subtraction meaningful and avoids relying on an unchecked address sum.
Once the candidate fits, its byte loop cannot run out of capacity halfway through.

Only after the entire spelling is stored does the routine update predecessor,
reading count and byte count. If the candidate does not fit, it returns the
previous prefix and writes none of the rejected candidate, even temporarily.
Writing bytes and then restoring their old values would violate this contract;
other observers can see those stores.

For the worked readings with byte capacity seven, four readings commit five
bytes. The final reading needs three, but only two positions remain. The return
is \texttt{\$v0}=4, \texttt{\$v1}=5, status one. Byte positions five and six
stay unchanged even though they lie inside the offered capacity. A capacity is
a maximum, not a promise that every available position will be used.

Empty input completes before testing capacity. Nonempty input at zero capacity
returns capacity status without reading the first reading. This ordering makes
the boundary cases part of the interface, not accidental consequences of loop
placement. The source shows the same distinction at the top of each iteration.

\section{Scan a candidate without advancing the prefix}

The decoder's arguments use the same registers, but their units change:
\texttt{\$a1} is byte length, and \texttt{\$a3} is word capacity. Results are
committed bytes in \texttt{\$v0}, committed readings in \texttt{\$v1}, and
status in \texttt{\$a3}. Besides complete and capacity, status two means
malformed bytes and status three means a reconstructed reading outside the
signed-word range. A nonempty stream with zero word capacity returns status one
before examining whether its next candidate is malformed.

The worked schedule keeps committed byte and reading counts in
\texttt{\$s5/\$s6}, and predecessor in \texttt{\$s4}. At the start of a
candidate, it copies the committed byte count to scan index \texttt{\$t0},
clears the unsigned accumulator, and begins at shift zero. These temporaries
may advance while the committed state stays fixed.

Before every byte load, check the scan index against logical input length.
An exhausted slice with a continued candidate is malformed, even if mapped
padding happens to contain a terminal byte. Mask out the payload, advance the
scan index, check the fifth-byte restrictions, then place the payload into the
unsigned accumulator. The payload shift increases by seven only when another
byte is required.

\Needspace{10\baselineskip}
\MARSCodeLabel{Accumulate one checked payload}
\begin{minted}{mars-idea}
dsllv $t6, $t5, $t2
or $t1, $t1, $t6
andi $t6, $t4, 128
beqzc $t6, dec_terminal
nop
\end{minted}
Here \texttt{\$t5} is the seven-bit payload, \texttt{\$t2} its bit position,
and \texttt{\$t4} the original byte. The preceding bounds and width checks
justify this fragment. Masking the continuation bit does not itself validate
that the byte exists or that its position is allowed.

Consider bytes \texttt{02 80 00}. The first byte commits reading one. The next
candidate scans both remaining bytes, obtaining unsigned zero, but fails the
shortest-spelling rule. At \texttt{dec\_bad}, scan index is three while committed
input is one and committed output is one. Returning three consumed bytes would
hide the rejected candidate from the caller. Returning no output would discard
the earlier valid reading.

\section{Validate the reconstructed reading before committing it}

After syntax succeeds, invert the unsigned mapping and add the full predecessor
in doubleword arithmetic. Only then check whether the candidate lies between
$-2147483648$ and 2147483647. A valid byte spelling need not describe a valid
reading in its current context. The syntactically valid five bytes
\texttt{ff ff ff ff 1f} map back to difference $-4294967296$; from predecessor
zero, that cannot be a signed-word reading and returns status three. A fifth
payload of 32 instead fails syntax and returns status two.

A word addition could wrap an invalid reconstructed value into the accepted
range. Keep the sum wide until both comparisons succeed. The decoder then
computes the full destination address, stores one word, updates predecessor,
copies the scan index to committed input, and increments committed output.
That ordering establishes the same prefix invariant as the encoder.

\Needspace{10\baselineskip}
\MARSCodeLabel{Commit one validated reading}
\begin{minted}{mars-idea}
sw $t6, 0($t7)
move $s4, $t6
move $s5, $t0
addiu $s6, $s6, 1
\end{minted}
The preceding address calculation and range checks are essential. An earlier
store would expose an invalid candidate. The function's return path restores
saved registers and stack pointer; inspect the working committed registers before
that restoration, or use the public results after return.

Both functions preserve all bits of saved registers, frame pointer and stack
pointer, return through the supplied continuation and perform no console I/O.
Word slices are four-byte aligned; byte slices may start anywhere. The existing
package permits an eight-byte-aligned incoming stack, while the worked caller
uses the book's sixteen-byte convention. Its 80-byte private frame accommodates
nine saved doublewords and preserves that alignment. The larger backing arenas
do not authorize reads outside selected input or writes beyond committed output.

\section{Compare representations on the same readings}

A useful comparison separates three decisions: fixed four-byte words, independent
ZigZag varints of the readings, and ZigZag varints of successive differences.
The middle choice benefits from small absolute values; the third benefits from
small changes. Neither makes every dataset smaller.

\begin{figure}[htbp]
  \centering
  \begin{tabular}{l r r r}
    \toprule
    Thirty-two readings           & Words & Independent varints & Delta varints \\
    \midrule
    Constant 1000                 & 128   & 64                  & 33            \\
    Ramp $-16$ through 15         & 128   & 32                  & 32            \\
    $10^9+1000i$, $i=0,\ldots,31$ & 128   & 160                 & 67            \\
    One 1000000 among zeros       & 128   & 34                  & 36            \\
    Alternating signed endpoints  & 128   & 160                 & 160           \\
    \bottomrule
  \end{tabular}
  \caption{Measured representation sizes in bytes, excluding framing. Independent mapping restarts from zero for each reading; delta mapping uses the previous reading within the block. The single spike is at index sixteen. All inputs and outputs are verified by the repeatable public measurement tool.}
  \label{fig:revision-delta-sizes}
\end{figure}


The ramp's values already fit in one mapped byte each, so delta coding cannot
improve its 32-byte total. The spike creates two large changes, one entering it
and one leaving it; independent coding has only one large absolute value.
Alternating signed endpoints need five bytes per item in either variable-length
scheme, exceeding the four-byte-word total. Choose a representation from the
actual distribution and required interfaces, not from one favorable example.

The repeatable \texttt{tools/measure\_revision\_delta.py} checks each complete
encoding, decodes its actual output in the same session, and records completed
instructions from function entry through return. For the worked source, steady
data takes 941 encoder and 1651 decoder instructions; alternating extremes take
2592 and 4000. Frame operations and validation are included; caller setup,
publication, assembly, session startup and host inspection are excluded. These
are simulator instruction counts, not hardware cycles or host elapsed time.
Different valid algorithms can have different work and the same bytes.

\section{Validate a stream independently}

Implement \texttt{scan\_delta} in \texttt{delta-scan.asm}. Pass a full byte
pointer in \texttt{\$a0}, byte length zero through 160 in \texttt{\$a1}, and
reading limit zero through 32 in \texttt{\$a2}. Return committed bytes and
readings in \texttt{\$v0/\$v1}, last committed signed reading in
\texttt{\$a2}, and the decoder's status in \texttt{\$a3}. If no reading has
committed, the last reading is zero. Apply exactly the same syntax, range and
precedence rules, but write no caller data, including temporary writes.

Explain which candidate state may change before acceptance and which state must
survive rejection. Use empty input, zero limit, a valid prefix followed by a
redundant zero, and a valid spelling whose reconstruction is out of range.
The default five-reading stream returns eight bytes, five readings, last reading
20000 and status zero. The lab supplies public checks that accept equivalent
implementations. Keep full-width pointers, preserved state and per-call reset
semantics; choose your own registers and algorithm.

\section*{Sources and laboratory connection}
The local format and item-atomic interface are specified in
\path{packages/023-delta-varint/}. The Protocol Buffers
\href{https://protobuf.dev/programming-guides/encoding/}{encoding guide}
describes base-128 varints and ZigZag; this lab does not implement Protobuf
messages. The \href{https://webassembly.github.io/spec/core/binary/values.html#integers}{WebAssembly integer grammar}
permits some redundant spellings that this lab rejects. Check the particular
format's rules when reusing a familiar byte loop. The revised lab investigates
committed prefixes, malformed candidates and the independent scanner; the
Base64 companion later applies bit grouping to a different representation task.
